\begin{figure*}[t]
\centering
\setcounter{subfigure}{0}
\begin{minipage}{170mm}
\ESPanel{.30\linewidth}{panels/consensus_method/e_matching.pdf}{Corresponding pairs at K}
\hfill
\ESPanel{.68\linewidth}{panels/consensus_method/b_support.pdf}{Count people whose layouts contain each corner}
\par\vspace{1pt}
\ESPanel{.63\linewidth}{panels/consensus_method/c_position.pdf}{Separate coordinate medians form a new upper--lower pair}
\hfill
\ESPanel{.35\linewidth}{panels/consensus_method/d_layout.pdf}{K in the complete collective layout}
\end{minipage}
\caption{Paired-corner structural consensus from 6 synthetic annotations. (a) Local walls locate the corresponding upper--lower pairs at K, with both endpoints used for matching; colors and symbols identify its 3 supporters. (b) Column layouts and row thumbnails identify who drew each corner. W has 6 supporters, K has exactly half and is retained, and E has only 2 and is omitted. A dot means the corner was not drawn. (c) K's upper elevation, unwrapped azimuth and lower elevation are sorted separately; orange boxes select each middle observation. Spacing shows rank, not angular distance. The upper and lower medians form a new pair with a shared median azimuth. (d) This pair defines corner K and its local wall height in the complete layout.}
\label{fig:consensus-steps}
\end{figure*}
